\SourceLeaf{015}\Article{10}
Obwohl die Beziehung zwischen \(\lambda\) und \(\mu\) von der Beschaffenheit der Funktion \(\varphi x\) abhängt, lassen sich doch einige allgemeine Ergebnisse feststellen. Wie diese Funktion auch beschaffen sei, sofern ihr Wert bei wachsendem Absolutwert von \(x\) stets abnimmt oder wenigstens nicht zunimmt, wird nämlich gewiss
\begin{quote}
\(\lambda\) kleiner oder wenigstens nicht größer als \(\mu\BareRoot{3}\) sein, wenn \(\mu\) kleiner als \(\tfrac23\) ist;

\(\lambda\) nicht größer als \(\frac{2}{3\BareRoot{(1-\mu)}}\) sein, wenn \(\mu\) größer als \(\tfrac23\) ist.
\end{quote}
Für \(\mu=\tfrac23\) fallen beide Grenzen zusammen: \(\lambda\) kann nämlich nicht größer als \(\BareRoot{\tfrac43}\) sein.

Um diesen bemerkenswerten Satz zu beweisen, bezeichnen wir mit \(y\) den Wert des Integrals \(\int\varphi z.\mathrm{d}z\) von \(z=-x\) bis \(z=+x\) erstreckt, sodass \(y\) die Wahrscheinlichkeit dafür ist, dass ein Fehler in den Grenzen \(-x\) und \(+x\) enthalten ist. Ferner setzen wir
\[
x=\psi y,\quad \mathrm{d}\psi y=\psi'y.\mathrm{d}y,\quad \mathrm{d}\psi'y=\psi''y.\mathrm{d}y
\]
Es wird also \(\psi0=0\) sein und ebenso
\[
\psi'y=\frac{1}{\varphi x+\varphi(-x)}
\]
weshalb nach Voraussetzung \(\psi'y\) von \(y=0\) bis \(y=1\) stets wachsen, wenigstens nirgends abnehmen wird; oder, was dasselbe ist, der Wert von \(\psi''y\) wird stets positiv oder wenigstens nicht negativ sein. Ferner haben wir \(\mathrm{d}.y\psi'y=\psi'y.\mathrm{d}y+y\psi''y.\mathrm{d}y\) und folglich
\[
y\psi'y-\psi y=\int y\psi''y.\mathrm{d}y
\]
wobei die Integration bei \(y=0\) beginnt. Der Wert des Ausdrucks \(y\psi'y-\psi y\) wird also stets eine positive, wenigstens nicht negative Größe sein und folglich
\[
1-\frac{\psi y}{y\psi'y}
\]
eine positive Größe kleiner als eins sein. Es sei \(f\) ihr Wert für \(y=\mu\), d. h. da \(\psi\mu=\lambda m\) gilt, sei
\[
f=1-\frac{\lambda m}{\mu\psi'\mu}\quad\text{oder}\quad\psi'\mu=\frac{\lambda m}{(1-f)\mu}
\]
Nach diesen Vorbereitungen betrachten wir die folgende Funktion von \(y\)
\[
\frac{\lambda m}{(1-f)\mu}(y-\mu f)
\]
die wir \(=Fy\) setzen wollen, sowie \(\mathrm{d}Fy=F'y.\mathrm{d}y\). Offensichtlich wird
\SourceLeaf{016}
\[
\begin{aligned}
F\mu&=\lambda m=\psi\mu\\[3pt]
F'\mu&=\frac{\lambda m}{(1-f)\mu}=\psi'\mu
\end{aligned}
\]
sein. Da nun \(\psi'y\) bei wachsendem \(y\) stetig zunimmt (wenigstens nicht abnimmt, was stets mitzudenken ist), \(F'y\) aber konstant ist, wird die Differenz \(\psi'y-F'y=\frac{\mathrm{d}(\psi y-Fy)}{\mathrm{d}y}\) für Werte von \(y\) größer als \(\mu\) positiv, für kleinere negativ sein. Hieraus folgt leicht, dass \(\psi y-Fy\) stets eine positive Größe ist und demnach \(\psi y\) stets absolut größer, wenigstens nicht kleiner als \(Fy\) sein wird, gewiss solange der Wert von \(Fy\) positiv ist, d. h. von \(y=\mu f\) bis \(y=1\). Daher wird der Wert des Integrals \(\int(Fy)^2\mathrm{d}y\) von \(y=\mu f\) bis \(y=1\) kleiner als der Wert des Integrals \(\int(\psi y)^2\mathrm{d}y\) zwischen denselben Grenzen sein und somit erst recht kleiner als der Wert dieses Integrals von \(y=0\) bis \(y=1\), der \(=mm\) ist. Der Wert des ersten Integrals ergibt sich aber als
\[
=\frac{\lambda\lambda mm(1-\mu f)^3}{3\mu\mu(1-f)^2}
\]
woraus wir schließen, dass \(\lambda\lambda\) kleiner als \(\frac{3\mu\mu(1-f)^2}{(1-\mu f)^3}\) ist, wobei \(f\) eine zwischen \(0\) und \(1\) liegende Größe ist. Nun nimmt der Wert des Bruchs \(\frac{3\mu\mu(1-f)^2}{(1-\mu f)^3}\), dessen Differential, wenn \(f\) als veränderliche Größe betrachtet wird, \(=\)
\[
-\frac{3\mu\mu(1-f)}{(1-\mu f)^4}.(2-3\mu+\mu f)\mathrm{d}f
\]
ist, stetig ab, während \(f\) vom Wert \(0\) zum Wert \(1\) übergeht, wenn \(\mu\) kleiner als \(\tfrac23\) ist; folglich ist der größtmögliche Wert derjenige, der \(f=0\) entspricht, nämlich \(=3\mu\mu\), sodass in diesem Fall \(\lambda\) gewiss kleiner oder nicht größer als \(\mu\BareRoot{3}\) wird. Q.~E.~P. Ist dagegen \(\mu\) größer als \(\tfrac23\), so wird der Wert jenes Bruchs für \(2-3\mu+\mu f=0\) am größten, d. h. für \(f=3-\tfrac2\mu\), woraus er sich als \(=\frac4{9(1-\mu)}\) ergibt, und somit ist in diesem Fall \(\lambda\) nicht größer als \(\frac2{3\BareRoot{(1-\mu)}}\). Q~E.~S.

So kann z. B. für \(\mu=\tfrac12\) gewiss \(\lambda\) nicht größer als \(\BareRoot{\tfrac34}\) sein, d. h. der wahrscheinliche Fehler kann die Grenze \(0{,}8660254m\) nicht überschreiten, der er im ersten Beispiel von Art. 9 gleich gefunden wurde. Ferner folgt aus unserem Satz leicht, dass \(\mu\) nicht kleiner als \(\lambda\BareRoot{\tfrac13}\) ist, solange \(\lambda\) kleiner als \(\BareRoot{\tfrac43}\) ist, dagegen \(\mu\) nicht kleiner als \(1-\frac4{9\lambda\lambda}\) für einen Wert von \(\lambda\) größer als \(\BareRoot{\tfrac43}\).

\Article{11}
Da mehrere im Folgenden zu behandelnde Aufgaben auch mit dem Wert des Integrals \(\int x^4\varphi x.\mathrm{d}x\) zusammenhängen, wird es der Mühe wert sein, ihn für einige beson\SourceLeaf{017}dere Fälle zu entwickeln. Wir bezeichnen den Wert dieses Integrals von \(x=-\infty\) bis \(x=+\infty\) erstreckt mit \(n^4\).

I. Für \(\varphi x=\tfrac1{2a}\), soweit \(x\) zwischen \(-a\) und \(+a\) enthalten ist, haben wir \(n^4=\tfrac15a^4=\tfrac95m^4\).

II. Im zweiten Fall von Art. 6, wo \(\varphi x=\frac{a\mp x}{aa}\) für Werte von \(x\) zwischen \(0\) und \(\pm a\) gilt, ergibt sich \(n^4=\tfrac1{15}a^4=\tfrac{12}{5}m^4\).

III. Im dritten Fall, wo
\[
\varphi x=\frac{e^{-\frac{xx}{hh}}}{a\BareRoot{\pi}}
\]
gilt, ergibt sich nach den Ausführungen in der oben zitierten Abhandlung \(n^4=\tfrac34h^4=3m^4\).

Im Übrigen lässt sich beweisen, dass der Wert von \(\frac{n^4}{m^4}\) gewiss nicht kleiner als \(\tfrac95\) ist, sofern die Voraussetzung des vorigen Artikels gilt.

\clearpage
